\section{External Consistency Check: CADE Cartel Convictions}
\label{sec:cade}

We check the FL classification against CADE (Conselho Administrativo
de Defesa Econ\^omica) cartel conviction records (2009--2019) as an
external consistency test. Of 65
CADE procurement cartel cases, 47 convicted firms match to BEC via
CNPJ. Among 2,735 FL firms, 193 (7.1\%) co-participate with at least
one CADE-convicted firm in the same tender-item. A permutation test
(1,000 iterations, stratified by participation quartile) yields a
permutation mean of 2.0\% (SD~$= 0.3$\%), implying that the observed
7.1\% is 3.5 times higher than expected ($p < 0.001$). That 92.9\%
of FL firms have no CADE overlap reflects the low base rate of cartel
prosecution: CADE convictions cover a small fraction of actual cartels
\citep{connor2007price}. The 7.1\% figure should be read as a lower
bound on the overlap rather than evidence that the remaining 92.9\% are
uninvolved.

\paragraph{What co-participation does and does not show.}
Sharing a tender-item with a convicted firm measures market proximity,
not cartel membership. An FL firm that co-participates may be a cover
bidder deployed by that cartel, an unwitting competitor in a
cartel-affected market, or a bystander. Importantly, once we condition
on participation count, the FL--CADE overlap rate becomes
indistinguishable from that of other always-losers ($p = 0.93$;
Section~\ref{sec:ground_truth}). The excess rate is therefore driven
largely by high participation volume rather than the specific FL
classification. This is not a deficiency but rather the screen's operating principle:
a participation-based triage tool flags firms that participate heavily
in suspicious markets. The gain is precisely that it works without
requiring the researcher to distinguish participation-driven proximity
from direct cartel involvement---a distinction that, with observational
data, may be impossible to draw cleanly. The permutation test
(3.5$\times$, stratified by participation quartile) confirms that the
excess is real, not mechanical. What the CADE check does not show is
that any given FL firm is a cartel member.

\paragraph{Dual pattern.}
Most CADE-convicted firms in BEC are \emph{winners}---the designated
recipients of cartel rents. But three convicted firms are themselves
classified as frequent losers (Table~\ref{tab:cade_fl_firms}). All
three participated across multiple item groups and years without ever
winning---textbook frequent-loser behavior.

\begin{table}[htbp]
\centering
\caption{CADE-convicted firms classified as frequent losers}
\label{tab:cade_fl_firms}
\small
\begin{tabular}{lllcc}
\toprule
Firm & Cartel & CADE process & Tenders & Wins \\
\midrule
Sol Tecnologia & Solar water heaters & 08012.001273/2010-24 & 84 & 0 \\
Nova Esperan\c{c}a Locadora & School transportation & 08700.005876/2019-85 & 97 & 0 \\
Jofran Com\'ercio & Trash bags (Op.\ Colludium) & 08700.005789/2015-02 & 65 & 0 \\
\bottomrule
\end{tabular}
\end{table}

\paragraph{Co-bidder analysis.}
FL firms appear in 5.6\% of CADE-firm tenders versus a 2.1\% baseline
(2.6$\times$). For individual convicted firms the rate is much
higher---up to 69\% for Mayfran (school transportation). Excluding
the three directly convicted FL firms and their 246 tender-items
leaves $\hat{\beta}$ unchanged at 0.064.

\paragraph{Beyond known cartels.}
Dropping all 31,447 CADE-involved tender-items and re-estimating the
baseline regression yields $\hat{\beta} = 0.062$ ($N =
1{,}622{,}954$), barely changed from the full-sample 0.064
(Table~\ref{tab:excl_cade}). The FL--price association does not
depend on tenders that CADE has already flagged, which means the
screen is at least picking up markets that fall outside the
enforcement record---whether because those markets harbor undetected
cartels or because FL presence correlates with other unobserved
features of high-price environments.

